\section{Estudio de Losas por Nivel}

\subsection{Determinación de Espesores de Losa}

Para el análisis, se seleccionaron las tres losas de mayores dimensiones de cada nivel. El procedimiento consistió en determinar la relación de aspecto $\epsilon$, interpolar el coeficiente de esbeltez $\phi$ (según las condiciones de apoyo indicadas en la normativa o apuntes de referencia) y aplicar la siguiente expresión para el cálculo del espesor mínimo:

\begin{equation}
    e \ge \frac{L_x \cdot \phi}{\lambda + r}
\end{equation}

Donde:
\begin{itemize}
    \item $L_x$: Luz menor del paño de losa [m].
    \item $\phi$: Coeficiente de esbeltez dependiente de las condiciones de borde.
    \item $\lambda$: Parámetro de diseño (35 para losas de piso tipo y 45 para losas de techumbre).
    \item $r$: Recubrimiento normativo considerado ($2 \text{ cm} = 0.02 \text{ m}$).
    \item $e$: Espesor calculado de la losa [m].
\end{itemize}

El espesor final de diseño para cada nivel se define como el valor máximo obtenido entre las tres losas analizadas, redondeado al centímetro superior por criterios constructivos.

\subsubsection{Piso 1}

Se analizaron las losas 101, 110 y 111. Los resultados obtenidos se presentan en la Tabla \ref{tab:esp_piso1}:

\begin{table}[H]
    \centering
    \begin{tabular}{lcccccc}
        \toprule
        \textbf{Losa} & \textbf{Tipo} & $L_x$ [m] & $L_y$ [m] & $\epsilon$ & $\phi$ & $e$ [cm] \\
        \midrule
        111 & 6  & 6.24 & 6.91 & 1.11 & 0.550 & 9.27 \\
        110 & 6  & 5.78 & 7.12 & 1.23 & 0.560 & 8.75 \\
        101 & 3a & 5.73 & 6.12 & 1.07 & 0.600 & 9.29 \\
        \bottomrule
    \end{tabular}
    \caption{Determinación de espesores para el primer nivel.}
    \label{tab:esp_piso1}
\end{table}

En consecuencia, se adopta un espesor final de \textbf{$e = 10 \text{ cm}$} para el primer piso.

\subsubsection{Piso 2}

Se analizaron las losas 201, 202 y 208. Los resultados se resumen a continuación:

\begin{table}[H]
    \centering
    \begin{tabular}{lcccccc}
        \toprule
        \textbf{Losa} & \textbf{Tipo} & $L_x$ [m] & $L_y$ [m] & $\epsilon$ & $\phi$ & $e$ [cm] \\
        \midrule
        201 & 6 & 6.24 & 6.91 & 1.11 & 0.550 & 9.27 \\
        202 & 6 & 6.18 & 8.43 & 1.36 & 0.567 & 9.47 \\
        208 & 6 & 5.73 & 8.75 & 1.52 & 0.591 & 9.14 \\
        \bottomrule
    \end{tabular}
    \caption{Determinación de espesores para el segundo nivel.}
    \label{tab:esp_piso2}
\end{table}

Se define un espesor de diseño de \textbf{$e = 10 \text{ cm}$} para el segundo piso.

\subsubsection{Pisos 3 al 21}

Para los niveles superiores (piso tipo), se analizaron las losas 301, 302 y 308:

\begin{table}[H]
    \centering
    \begin{tabular}{lcccccc}
        \toprule
        \textbf{Losa} & \textbf{Tipo} & $L_x$ [m] & $L_y$ [m] & $\epsilon$ & $\phi$ & $e$ [cm] \\
        \midrule
        301 & 6 & 5.73 & 8.75 & 1.52 & 0.591 & 9.14 \\
        302 & 6 & 6.18 & 8.43 & 1.36 & 0.567 & 9.47 \\
        308 & 6 & 6.21 & 8.98 & 1.45 & 0.575 & 9.65 \\
        \bottomrule
    \end{tabular}
    \caption{Determinación de espesores para pisos tipo (3 al 21).}
    \label{tab:esp_pisotipo}
\end{table}

Se adopta un espesor de diseño de \textbf{$e = 10 \text{ cm}$} para la planta tipo.

\subsubsection{Subterráneo}

En el nivel de subterráneo se consideraron las losas 008, 035 y 036. Se destaca un incremento significativo en el espesor debido a las condiciones de apoyo de la losa 036:

\begin{table}[H]
    \centering
    \begin{tabular}{lcccccc}
        \toprule
        \textbf{Losa} & \textbf{Tipo} & $L_x$ [m] & $L_y$ [m] & $\epsilon$ & $\phi$ & $e$ [cm] \\
        \midrule
        008 & 6  & 6.24 & 6.91 & 1.11 & 0.550 & 9.27 \\
        035 & 6  & 5.78 & 7.12 & 1.23 & 0.560 & 8.75 \\
        036 & 2b & 5.60 & 8.15 & 1.45 & 0.945 & 14.30 \\
        \bottomrule
    \end{tabular}
    \caption{Determinación de espesores para el nivel subterráneo.}
    \label{tab:esp_subt}
\end{table}

Debido a los requerimientos de la losa 036, el espesor final para el subterráneo se define en \textbf{$e = 15 \text{ cm}$}.